\section{Discussion}

The paper studies frequency-report scoring rules as inferential devices.
The report \(r\) is not the belief vector.
It is evidence about \(p\), and the scoring rule determines how informative that evidence is.
The relevant chain is therefore
\[
S
\rightarrow
R_S(p)
\rightarrow
P_S(r)
\rightarrow
\text{bounds on beliefs and functionals}.
\]

The results suggest three practical recommendations.
First, use discrete-metric frequency scoring when the priority is transparency and fixed-prize robustness.
It has closed-form coordinate bounds and is robust to risk aversion, but exact-match payment probabilities can be low away from concentrated reports.

Second, use quadratic-distance frequency scoring when the priority is transparent belief and mean inference.
It produces a projection rule, a linear belief region, closed-form coordinate bounds, and linear-program bounds for any linear functional.
The design exercise shows that quadratic distance is especially attractive for worst-coordinate uncertainty and for mean-oriented targets in sparse-belief environments.

Third, use Manhattan distance when median-type frequency reports are substantively desired or when it is useful as a comparison rule.
Its binary case is clean, and its multi-category inverse region is explicit, but the coordinate and mean bounds require threshold computation rather than a simple closed form.

The paper should not be read as a general theory of all discrete scoring rules.
Hamming and Chebyshev distance are discussed as additional computable rules, but the primary analytic-bound results are the quadratic-distance and discrete-metric cases.
The intended contribution is narrower: a finite-sample design guide for researchers who want to interpret observed frequency reports as bounds on latent beliefs and target functionals.

Important open questions remain.
The design exercise should be refined before empirical use, especially with alternative outcome vectors, a final choice of paper tables or figures, and robustness checks for tie-breaking conventions.
The cognitive advantage of frequency reports and distance payments should also be tested experimentally rather than assumed.
Those extensions are separate from the main theoretical point: different frequency-report scoring rules induce different inverse belief regions, and those regions determine the informational efficiency of the mechanism.
